% Aula 10 --- crescer com Gini, podar com o erro: a poda remove os ramos que não mudam nenhuma decisão.
Rafael é analista numa operadora de planos de saúde e cresceu uma árvore de classificação,
com o índice de Gini, para prever quais beneficiários serão internados no próximo ano
($Y=1$). Agora vai podá-la. Para uma região $R$ com $n_R$ observações de treino, sejam
$n_{R,c}$ o número delas na classe $c$ e $\widehat p_{R,c}=n_{R,c}/n_R$; o índice de Gini de
$R$ é $G(R)=\sum_c\widehat p_{R,c}(1-\widehat p_{R,c})$. A folha $R$ prediz a classe
majoritária, e $E(R)=1-\max_c\widehat p_{R,c}$ é o seu erro de classificação, de modo
que $n_RE(R)=n_R-\max_cn_{R,c}$ é o número de erros de treino na folha. Na poda por
custo-complexidade com o erro de classificação, escolhe-se, entre as subárvores $T$ da árvore
crescida, a que minimiza
\[
  C_\alpha(T)=\sum_{m=1}^{\abs{T}}n_mE(R_m)+\alpha\abs{T},
\]
em que $R_1,\dots,R_{\abs{T}}$ são as folhas de $T$, $n_m=n_{R_m}$ e $\abs{T}$ é o número de
folhas; o valor de $\alpha>0$ é escolhido por validação cruzada.
\begin{itens}
  \item \pts{1,3} Seja $R$ uma região dividida em regiões disjuntas $R_1,\dots,R_L$. Mostre
        que $\sum_ln_{R_l}E(R_l)\le n_RE(R)$, com qualquer número de classes, e que vale a
        igualdade quando uma mesma classe $c^*$ é majoritária em todas as $R_l$; nesse caso,
        $c^*$ é majoritária também em $R$.
  \item \pts{1,6} Mostre que, para todo $\alpha>0$, nenhuma subárvore que minimize $C_\alpha$
        tem um nó interno cujas folhas descendentes predizem todas a mesma classe.
  \item \pts{1,1} Na árvore de Rafael, um nó $R$ com $400$ beneficiários, dos quais $100$
        foram internados, foi dividido em dois filhos de $200$, que são folhas da árvore
        crescida: $R_1$, com $20$ internados, e $R_2$, com $80$. Calcule a redução do índice
        de Gini ponderado,
        $G(R)-\big[n_1G(R_1)+n_2G(R_2)\big]/n_R$, e a redução do erro de classificação
        ponderado. Por que o crescimento fez essa divisão, e o que a poda fará com ela?
\end{itens}

\begin{solucao}
\textbf{(a)} As regiões $R_l$ particionam $R$, e por isso $n_{R,c}=\sum_ln_{R_l,c}$ para
cada classe $c$. O máximo de uma soma não passa da soma dos máximos:
\[
  \max_cn_{R,c}=\max_c\sum_ln_{R_l,c}\ \le\ \sum_l\max_cn_{R_l,c} .
\]
Logo
\[
  \sum_ln_{R_l}E(R_l)=\sum_l\Big(n_{R_l}-\max_cn_{R_l,c}\Big)=n_R-\sum_l\max_cn_{R_l,c}
  \ \le\ n_R-\max_cn_{R,c}=n_RE(R).
\]
Se $c^*$ é majoritária em todas as $R_l$, então
$\sum_l\max_cn_{R_l,c}=\sum_ln_{R_l,c^*}=n_{R,c^*}\le\max_cn_{R,c}$; com a desigualdade
anterior, os três números são iguais. A igualdade vale, e $n_{R,c^*}=\max_cn_{R,c}$: $c^*$ é
majoritária em $R$. Dividir uma região nunca aumenta o número de erros de treino, e só o
diminui quando alguma das partes muda a classe predita.

\smallskip
\textbf{(b)} Suponha que $T$ minimize $C_\alpha$ e tenha um nó interno $t$ cujas
$L\ge2$ folhas descendentes predizem todas a classe $c^*$. Seja $T'$ a subárvore obtida
trocando o ramo que sai de $t$ por uma única folha, a região de $t$. Pelo item (a), essa folha
tem $c^*$ como classe majoritária e o mesmo número de erros de treino que as $L$ folhas
juntas; as outras folhas não mudam. As duas subárvores cometem, então, os mesmos erros, e
$\abs{T'}=\abs{T}-(L-1)$, de modo que
\[
  C_\alpha(T')=C_\alpha(T)-\alpha(L-1)<C_\alpha(T),
\]
porque $\alpha>0$ e $L\ge2$. Isso contradiz a escolha de $T$. Qualquer que seja $\alpha>0$, a
poda pelo erro de classificação remove todos os ramos que não mudam nenhuma decisão.

\smallskip
\textbf{(c)} As proporções de internados são $0{,}25$ no nó e $0{,}1$ e $0{,}4$ nos filhos. Com
duas classes, as duas parcelas de $G$ são iguais, e $G=2\widehat p(1-\widehat p)$, com
$\widehat p$ a proporção de internados:
\[
  G(R)=2\cdot0{,}25\cdot0{,}75=0{,}375,
  \qquad
  \tfrac12\,G(R_1)+\tfrac12\,G(R_2)=\tfrac12\cdot0{,}18+\tfrac12\cdot0{,}48=0{,}33 ,
\]
uma redução de $0{,}045$. A classe majoritária é ``não internado'' no nó ($300$ contra $100$) e
nos dois filhos ($180$ contra $20$ e $120$ contra $80$). Os erros são $100$ no nó e
$20+80=100$ nos filhos: a redução do erro é zero, como o item (a) previa. O Gini é
estritamente côncavo e registra o ganho de pureza de separar um grupo de risco baixo
($10\%$ de internações) de um de risco alto ($40\%$); por isso o crescimento fez a divisão.
Como os filhos são folhas da árvore crescida, numa subárvore que mantenha a divisão eles são
as únicas folhas descendentes de $R$, e os dois predizem a mesma classe. Pelo item (b), a poda
pelo erro remove a divisão, qualquer que seja $\alpha>0$.
\end{solucao}
